% 02-metodo.tex — Sección 2: método de medición y verificación del baseline.
% Responde las preguntas 1, 2 y 5 del revisor (diapo 11/13) y declara las
% decisiones de SPEC §5 en vez de esconderlas (criterio de PLAN F9).

\section{Método de medición y verificación del baseline}
\label{sec:metodo}

\subsection{El conjunto de datos es el mismo para los dos motores}
\label{sec:baseline}

Los \num{500000} documentos se generan con
\texttt{generator/synthetic\_data\_generator.py} con \textbf{semilla fija 42}, de
modo que el archivo es reproducible byte a byte. Su huella SHA-256 empieza por
\texttt{0cd7fadb} y termina en \texttt{0414}, y queda registrada en
\texttt{results/dataset\_stats.json} junto con las distribuciones obtenidas.

Que los dos motores contengan lo mismo no se supone: se verificó en tres
direcciones (archivo~$\to$~motor, motor~$\to$~archivo y motor~$\to$~motor) con
\texttt{scripts/verify\_gate.py}, comparando los documentos exportados de cada
motor contra la fuente en \emph{streaming}. Dos decisiones del esquema son las
que hacen que la comparación sea válida:

\begin{itemize}
  \item \texttt{payment.final\_total} se almacena como \textbf{número} JSON, no
        como cadena. Si fuera cadena, MongoDB ordenaría lexicográficamente
        (\texttt{"9.5" > "100.0"}) y PostgreSQL numéricamente: los Top-1000 no
        coincidirían y la comparación quedaría invalidada \emph{en silencio}.
  \item \texttt{timestamp} se almacena como cadena ISO-8601 en los dos motores.
        Como \texttt{Date} en MongoDB, el \texttt{\$gte} contra un literal de
        texto no haría \emph{match}.
\end{itemize}

\subsection{Las consultas se leen del disco, no están dentro del código}

Las cuatro consultas (\texttt{Q1} y \texttt{Q2}, para cada motor) se copiaron
\textbf{literalmente} del manual a \texttt{queries/} y el runner las lee de ahí en
cada ejecución, con una comprobación de integridad al arrancar. El revisor puede
compararlas con el PDF sin leer código.

Su equivalencia se verificó antes de medir: \texttt{scripts/check\_equivalence.py}
calcula el resultado de referencia en Python con aritmética exacta
(\texttt{Decimal}) y lo compara con lo que devuelve cada motor. Coinciden en todas
las variantes. El orden \emph{dentro} del Top-1000 de \texttt{Q2} no es
determinista —hay valores repetidos y la consulta no tiene desempate—, así que se
compara como conjunto.

\subsection{Qué es una repetición}

\begin{enumerate}
  \item Se baja el motor que estuviera arriba. \textbf{Solo un motor corre a la
        vez}: con los dos arriba competirían por la caché de páginas de la VM y la
        comparación dejaría de ser válida.
  \item Se limpia la caché y \textbf{se mide el efecto} (cuánto bajó
        \texttt{buff/cache}), no solo el código de salida del comando. Entre 0.7 y
        1.8~GB liberados por repetición, registrados en el archivo crudo.
  \item Se levanta el motor con el perfil de la celda y se espera su
        \emph{healthcheck}.
  \item \textbf{Se verifican los límites leyendo el cgroup} y se comparan con los
        esperados del perfil. Si no coinciden, la repetición se aborta sin
        escribir fila.
  \item Ventanas de \textbf{60~s de calentamiento, 300~s de medición y 60~s de
        enfriamiento}, como fija el manual. Los clientes trabajan en bucle cerrado,
        sin pausa entre consultas.
\end{enumerate}

La concurrencia 50 se reparte en \textbf{10 procesos de 5 hilos}, cada hilo con su
propia conexión. Se verificó que el motor ve los 50 clientes simultáneos
igual que con 50 procesos, y se eligió este reparto porque con 50 procesos el
runner ocupaba unos 2~GB y empujaba la VM a \emph{swap}, frente a unos 460~MB con
10 procesos. La decisión se apoya en el P50, el throughput, la concurrencia vista
por el motor, las conexiones y la memoria; \textbf{no} en el P95, que quedó
contaminado por el arranque en frío.

\subsection{Cómo se resumen las repeticiones (pregunta 5 del revisor)}
\label{sec:resumen}

\textbf{Tres repeticiones por celda}, y las tablas de este reporte dan la
\textbf{media} de las tres. Se eligió la media y no la mediana porque con tres
valores la mediana es simplemente el del medio y descarta los otros dos. Para que
el promedio no esconda dispersión, cada celda se acompaña de su coeficiente de
variación, y las celdas cuya dispersión supera el 15~\% o que muestran una
tendencia monótona se señalan una por una en la sección~\ref{sec:resultados}.

\subsection{Cómo se calculan los percentiles (pregunta 2 del revisor)}

Se reportan \textbf{P50, P95 y P99}, no promedios de latencia. Se calculan con
interpolación lineal (\texttt{statistics.quantiles} con
\texttt{method="inclusive"}, el mismo método que \texttt{numpy} usa por defecto)
sobre las latencias de las consultas \textbf{completadas dentro de la ventana de
medición}. Las latencias crudas de cada consulta se conservan en
\texttt{results/raw/}, así que los percentiles se pueden recalcular con otro
método sin volver a medir.

\paragraph{Los \emph{timeouts} son datos censurados, no ceros.}
Una consulta que agota el límite de 300~s no se descarta ni se cuenta como un
valor: es un dato \textbf{censurado por la derecha}, del que solo se sabe que su
latencia es mayor que el límite. Por eso, cuando la fracción de \emph{timeouts}
dentro de la ventana alcanza el 1~\%, la fila se marca con
\texttt{status = error} y \textbf{cada percentil se declara válido o no}: el P99
solo vale si los \emph{timeouts} son menos del 1~\%, el P95 menos del 5~\% y el
P50 menos del 50~\%. Esa validez va fila por fila en \texttt{results/raw/}. En
este experimento \cifraCeldasCensuradas{} celdas quedaron con percentiles
invalidados, y se marcan en todas las tablas.

\paragraph{Celdas con pocas muestras.}
Con 300~s de ventana y consultas de 20~s, una celda puede completar 15
consultas. Se marcan como \textbf{baja densidad muestral} (menos de 20
completadas) y su P99 se lee con esa reserva, en vez de presentarlo como si
viniera de cientos de muestras.

\subsection{De dónde sale cada columna del CSV}

Las 14 columnas se miden directamente, y \textbf{ninguna se rellena con un cero
inventado}: si una lectura falta, la celda queda vacía y la fila sale con
\texttt{status} \texttt{error} u \texttt{oom}.

\begin{table}[htbp]
  \centering
  \caption{Origen de las columnas que no son identificadores.}
  \label{tab:origen-columnas}
  \small
  \begin{tabularx}{\linewidth}{lX}
    \toprule
    \texttt{throughput\_qps} & consultas completadas en la ventana $\div$ 300~s \\
    \texttt{latency\_p50/p95/p99\_ms} & percentiles de las completadas en la ventana \\
    \texttt{cpu\_throttled\_seconds} & $\Delta$ \texttt{throttled\_usec} de
      \texttt{cpu.stat} del cgroup $\div$ \num{1e6} \\
    \texttt{max\_rss\_mb} & máximo de \texttt{anon} de \texttt{memory.stat},
      muestreado cada segundo durante la ventana \\
    \texttt{disk\_spill} & $\Delta$ \texttt{temp\_bytes} (PostgreSQL) o
      contadores de \emph{spill} de \texttt{serverStatus} (MongoDB) \\
    \texttt{status} & \texttt{oom} > \texttt{error} > \texttt{spill} > \texttt{ok} \\
    \bottomrule
  \end{tabularx}
\end{table}

El cgroup se lee \textbf{desde fuera del contenedor}, en
\texttt{/sys/fs/cgroup/docker/}, y no con \texttt{docker exec}: se midió que cada
lectura por \texttt{exec} cuesta unos 250~ms y gasta CPU \emph{del motor}, y
muestrear así cada segundo en el perfil C se comería un 8~\% de su cuota. Es
decir, mediríamos la medición.

\paragraph{Jerarquía de fuentes.}
Las columnas del CSV salen del cgroup y de los motores. Prometheus, cAdvisor y
node\_exporter (con \texttt{scrape\_interval: 1s}) alimentan las curvas y sirven
de control cruzado, pero \textbf{no} producen ninguna columna. Se verificó que
las dos fuentes miden lo mismo: \texttt{container\_memory\_rss} de cAdvisor es
exactamente \texttt{anon} del cgroup —comprobado con el contenedor congelado, donde
ningún proceso puede alterar la memoria— y el throttling de las dos fuentes
difiere menos del 1~\% bajo carga.

\subsection{Detección de suspensión del equipo}

Si el equipo se suspende durante una ventana, la VM se congela: las consultas en
vuelo se detienen y el reloj de pared avanza sin que corra nada. Esa ventana ya no
mide lo que dice medir. El runner lo detecta comparando el reloj de pared con el
monótono durante la ventana y marca la repetición como \texttt{error}. En esta
matriz \textbf{no ocurrió ninguna vez}, y el número de repeticiones descartadas
por cualquier motivo es \cifraDescartadas{} (sección~\ref{sec:integridad}).
